\documentclass[10pt,a4paper]{amsart}
\usepackage{amsmath,amssymb,amsthm,graphicx}
% Retained preamble for the Stacks Project chapter {CH} (shared project preamble).
% The chapter file pinned-original-source.tex only carries \input{{preamble}} on its first line,
% so no preamble text lives inside the chapter itself.  The Stacks Project uses one single shared
% preamble for every chapter: the snapshot stages it three times (sources/stacks-chow/preamble.tex,
% sources/stacks-etale-cohomology/preamble.tex and focused-seeds/AGP/sources/stacks-snapshot-20261011/
% preamble.tex) and those copies are identical in size.  The block below is copied mechanically,
% line for line, from sources/stacks-chow/preamble.tex lines 14-18 and 154-217: the Xy-pic setup and
% the native theorem declarations together with the mathematical macros.  Nothing was added, removed
% or redefined; the wrapper class and the amsmath/amssymb/amsthm/graphicx base packages are supplied
% by the assembler.
% For commutative diagrams we use Xy-pic
% Chapter: stacks-cohomology.tex -- Cohomology of Sheaves.
\usepackage[all]{xy}

% We use 2cell for 2-commutative diagrams.
\xyoption{2cell}

% Theorem environments.
%
\theoremstyle{plain}
\newtheorem{theorem}[subsection]{Theorem}
\newtheorem{proposition}[subsection]{Proposition}
\newtheorem{lemma}[subsection]{Lemma}

\theoremstyle{definition}
\newtheorem{definition}[subsection]{Definition}
\newtheorem{example}[subsection]{Example}
\newtheorem{exercise}[subsection]{Exercise}
\newtheorem{situation}[subsection]{Situation}

\theoremstyle{remark}
\newtheorem{remark}[subsection]{Remark}
\newtheorem{remarks}[subsection]{Remarks}

\numberwithin{equation}{subsection}

% Macros
%
\def\lim{\mathop{\mathrm{lim}}\nolimits}
\def\colim{\mathop{\mathrm{colim}}\nolimits}
\def\Spec{\mathop{\mathrm{Spec}}}
\def\Hom{\mathop{\mathrm{Hom}}\nolimits}
\def\Ext{\mathop{\mathrm{Ext}}\nolimits}
\def\SheafHom{\mathop{\mathcal{H}\!\mathit{om}}\nolimits}
\def\SheafExt{\mathop{\mathcal{E}\!\mathit{xt}}\nolimits}
\def\Sch{\mathit{Sch}}
\def\Mor{\mathop{\mathrm{Mor}}\nolimits}
\def\Ob{\mathop{\mathrm{Ob}}\nolimits}
\def\Sh{\mathop{\mathit{Sh}}\nolimits}
\def\NL{\mathop{N\!L}\nolimits}
\def\CH{\mathop{\mathrm{CH}}\nolimits}
\def\proetale{{pro\text{-}\acute{e}tale}}
\def\etale{{\acute{e}tale}}
\def\QCoh{\mathit{QCoh}}
\def\Ker{\mathop{\mathrm{Ker}}}
\def\Im{\mathop{\mathrm{Im}}}
\def\Coker{\mathop{\mathrm{Coker}}}
\def\Coim{\mathop{\mathrm{Coim}}}

% Boxtimes
%
\DeclareMathSymbol{\boxtimes}{\mathbin}{AMSa}{"02}

%
% Macros for moduli stacks/spaces
%
\def\QCohstack{\mathcal{QC}\!\mathit{oh}}
\def\Cohstack{\mathcal{C}\!\mathit{oh}}
\def\Spacesstack{\mathcal{S}\!\mathit{paces}}
\def\Quotfunctor{\mathrm{Quot}}
\def\Hilbfunctor{\mathrm{Hilb}}
\def\Curvesstack{\mathcal{C}\!\mathit{urves}}
\def\Polarizedstack{\mathcal{P}\!\mathit{olarized}}
\def\Complexesstack{\mathcal{C}\!\mathit{omplexes}}
% \Pic is the operator that assigns to X its picard group, usage \Pic(X)
% \Picardstack_{X/B} denotes the Picard stack of X over B
% \Picardfunctor_{X/B} denotes the Picard functor of X over B
\def\Pic{\mathop{\mathrm{Pic}}\nolimits}
\def\Picardstack{\mathcal{P}\!\mathit{ic}}
\def\Picardfunctor{\mathrm{Pic}}
\def\Deformationcategory{\mathcal{D}\!\mathit{ef}}
\usepackage{cleveref}
\crefname{originallemma}{original lemma}{original lemmas}
\Crefname{originallemma}{Original lemma}{Original lemmas}
\makeatletter
\newlabel{lemma-alternating-cech-complex-complex}{{08C1}{source}{}{sourcepointer-1}{}}
\newlabel{lemma-alternating-cech-complex-complex@cref}{{[originallemma][0][]08C1}{[1][0][]source}{}{}{}}
\makeatother
\crefname{originallemma}{original lemma}{original lemmas}
\Crefname{originallemma}{Original lemma}{Original lemmas}
\makeatletter
\newlabel{lemma-alternating-usual}{{01FM}{source}{}{sourcepointer-2}{}}
\newlabel{lemma-alternating-usual@cref}{{[originallemma][0][]01FM}{[1][0][]source}{}{}{}}
\makeatother
\crefname{originalreference}{original reference}{original references}
\Crefname{originalreference}{Original reference}{Original references}
\makeatletter
\newlabel{homology-lemma-ss-double-complex}{{0130}{source}{}{sourcepointer-3}{}}
\newlabel{homology-lemma-ss-double-complex@cref}{{[originalreference][0][]0130}{[1][0][]source}{}{}{}}
\makeatother
\crefname{originalreference}{original reference}{original references}
\Crefname{originalreference}{Original reference}{Original references}
\makeatletter
\newlabel{homology-lemma-first-quadrant-ss}{{0132}{source}{}{sourcepointer-4}{}}
\newlabel{homology-lemma-first-quadrant-ss@cref}{{[originalreference][0][]0132}{[1][0][]source}{}{}{}}
\makeatother
\crefname{originallemma}{original lemma}{original lemmas}
\Crefname{originallemma}{Original lemma}{Original lemmas}
\makeatletter
\newlabel{lemma-cech-complex-complex-computes}{{0FLH}{source}{}{sourcepointer-5}{}}
\newlabel{lemma-cech-complex-complex-computes@cref}{{[originallemma][0][]0FLH}{[1][0][]source}{}{}{}}
\makeatother
\crefname{originallemma}{original lemma}{original lemmas}
\Crefname{originallemma}{Original lemma}{Original lemmas}
\makeatletter
\newlabel{lemma-K-injective}{{071B}{source}{}{sourcepointer-6}{}}
\newlabel{lemma-K-injective@cref}{{[originallemma][0][]071B}{[1][0][]source}{}{}{}}
\makeatother
\crefname{originalreference}{original reference}{original references}
\Crefname{originalreference}{Original reference}{Original references}
\makeatletter
\newlabel{homology-lemma-apply-Mittag-Leffler-again}{{070E}{source}{}{sourcepointer-7}{}}
\newlabel{homology-lemma-apply-Mittag-Leffler-again@cref}{{[originalreference][0][]070E}{[1][0][]source}{}{}{}}
\makeatother
\crefname{originalreference}{original reference}{original references}
\Crefname{originalreference}{Original reference}{Original references}
\makeatletter
\newlabel{more-algebra-lemma-apply-Mittag-Leffler-again}{{0918}{source}{}{sourcepointer-8}{}}
\newlabel{more-algebra-lemma-apply-Mittag-Leffler-again@cref}{{[originalreference][0][]0918}{[1][0][]source}{}{}{}}
\makeatother
\title{The alternating Cech complex computes derived global sections of an unbounded complex of modules under vanishing hypotheses on a basis of opens}
\author{The Stacks Project Authors}
\date{stacks-project:cohomology.tex:local-snapshot-20261011; GNU FDL 1.2 or later}
\begin{document}
\maketitle
\noindent\textit{Editorial note.} All delivered mathematical text is the author's own, copied verbatim from the pinned Stacks Project chapter file cohomology.tex; no statement or proof step was rewritten, reordered or removed.\par\medskip
\noindent\textit{Editorial note.} The counted span cites eight results living in earlier chapters; each is entered in assembly.external\_references as a navigation pointer with its canonical tag, and none of them replaces a load-bearing step of this proof (the unbounded limit argument is fully inside the span).\par\medskip
\noindent\textit{Editorial note.} The proof cites no bibliography item inside the counted span, so the card carries no bibliography excerpt.\par\medskip
\noindent\textit{Allowed standard prerequisites (editorial).}\begin{itemize}
\item Comparison of the usual Cech complex with derived global sections for bounded below complexes. Stacks Project, Cohomology of Sheaves, Lemma lemma-cech-complex-complex-computes (tag 0FLH). Supplies the bounded below case once the alternating Cech complex is shown quasi-isomorphic to the usual one.
\item Spectral sequences of a double complex and their convergence. Stacks Project, Homology, Lemmas homology-lemma-ss-double-complex (tag 0130) and homology-lemma-first-quadrant-ss (tag 0132). Justify the first-page comparison and the convergence step in the bounded below case.
\item Mittag-Leffler criterion for the cohomology of a limit of complexes. Stacks Project, Homology, Lemma homology-lemma-apply-Mittag-Leffler-again (tag 070E); also More on Algebra, Lemma more-algebra-lemma-apply-Mittag-Leffler-again (tag 0918). Lets the eventual constancy of the left-hand cohomology be transferred to the limit on the right-hand side.
\end{itemize}\medskip
\ifdefined\hypertarget\hypertarget{sourcepointer-1}{}\fi
\noindent\textit{Source reference.} Original lemma 08C1: sources/stacks-chow/tags.txt, source lines 10586--10586. Cross-chapter standard prerequisite cited by the author inside the counted span; it is used as a pointer and its own proof is not part of this item.\par\medskip
\ifdefined\hypertarget\hypertarget{sourcepointer-2}{}\fi
\noindent\textit{Source reference.} Original lemma 01FM: sources/stacks-chow/tags.txt, source lines 1875--1875. Cross-chapter standard prerequisite cited by the author inside the counted span; it is used as a pointer and its own proof is not part of this item.\par\medskip
\ifdefined\hypertarget\hypertarget{sourcepointer-3}{}\fi
\noindent\textit{Source reference.} Original reference 0130: sources/stacks-chow/tags.txt, source lines 1418--1418. Cross-chapter standard prerequisite cited by the author inside the counted span; it is used as a pointer and its own proof is not part of this item.\par\medskip
\ifdefined\hypertarget\hypertarget{sourcepointer-4}{}\fi
\noindent\textit{Source reference.} Original reference 0132: sources/stacks-chow/tags.txt, source lines 1420--1420. Cross-chapter standard prerequisite cited by the author inside the counted span; it is used as a pointer and its own proof is not part of this item.\par\medskip
\ifdefined\hypertarget\hypertarget{sourcepointer-5}{}\fi
\noindent\textit{Source reference.} Original lemma 0FLH: sources/stacks-chow/tags.txt, source lines 19497--19497. Cross-chapter standard prerequisite cited by the author inside the counted span; it is used as a pointer and its own proof is not part of this item.\par\medskip
\ifdefined\hypertarget\hypertarget{sourcepointer-6}{}\fi
\noindent\textit{Source reference.} Original lemma 071B: sources/stacks-chow/tags.txt, source lines 8985--8985. Cross-chapter standard prerequisite cited by the author inside the counted span; it is used as a pointer and its own proof is not part of this item.\par\medskip
\ifdefined\hypertarget\hypertarget{sourcepointer-7}{}\fi
\noindent\textit{Source reference.} Original reference 070E: sources/stacks-chow/tags.txt, source lines 8953--8953. Cross-chapter standard prerequisite cited by the author inside the counted span; it is used as a pointer and its own proof is not part of this item.\par\medskip
\ifdefined\hypertarget\hypertarget{sourcepointer-8}{}\fi
\noindent\textit{Source reference.} Original reference 0918: sources/stacks-chow/tags.txt, source lines 11436--11436. Cross-chapter standard prerequisite cited by the author inside the counted span; it is used as a pointer and its own proof is not part of this item.\par\medskip
\begin{lemma}
\label{lemma-alternating-cech-complex-complex-ss}
Let $(X, \mathcal{O}_X)$ be a ringed space. Let
$\mathcal{U} : X = \bigcup_{i \in I} U_i$ be a finite open covering. Let
$\mathcal{F}^\bullet$ be a complex of $\mathcal{O}_X$-modules.
Let $\mathcal{B}$ be a set of open subsets of $X$. Assume
\begin{enumerate}
\item every open in $X$ has a covering whose members are
elements of $\mathcal{B}$,
\item we have $U_{i_0\ldots i_p} \in \mathcal{B}$ for all
$i_0, \ldots, i_p \in I$,
\item for every $U \in \mathcal{B}$ and $p > 0$ we have
\begin{enumerate}
\item $H^p(U, \mathcal{F}^q) = 0$,
\item $H^p(U, \Coker(\mathcal{F}^{q - 1} \to \mathcal{F}^q)) = 0$, and
\item $H^p(U, H^q(\mathcal{F})) = 0$.
\end{enumerate}
\end{enumerate}
Then the map
$$
\text{Tot}(\check{\mathcal{C}}^\bullet_{alt}(\mathcal{U}, \mathcal{F}^\bullet))
\longrightarrow
R\Gamma(X, \mathcal{F}^\bullet)
$$
of Lemma \ref{lemma-alternating-cech-complex-complex}
is an isomorphism in $D(\textit{Ab})$.
\end{lemma}

\begin{proof}
First assume $\mathcal{F}^\bullet$ is bounded below. In this case the map
$$
\text{Tot}(\check{\mathcal{C}}^\bullet_{alt}(\mathcal{U}, \mathcal{F}^\bullet))
\longrightarrow
\text{Tot}(\check{\mathcal{C}}^\bullet(\mathcal{U}, \mathcal{F}^\bullet))
$$
is a quasi-isomorphism by Lemma \ref{lemma-alternating-usual}.
Namely, the map of double complexes
$\check{\mathcal{C}}^\bullet_{alt}(\mathcal{U}, \mathcal{F}^\bullet) \to
\check{\mathcal{C}}^\bullet(\mathcal{U}, \mathcal{F}^\bullet)$
induces an isomorphism between the first pages of the second spectral sequences
associated to these complexes
(by Homology, Lemma \ref{homology-lemma-ss-double-complex})
and these spectral sequences converge
(Homology, Lemma \ref{homology-lemma-first-quadrant-ss}).
Thus the conclusion in this case by
Lemma \ref{lemma-cech-complex-complex-computes} and assumption (3)(a).

\medskip\noindent
In general, by assumption (3)(c) we may choose a resolution
$\mathcal{F}^\bullet \to \mathcal{I}^\bullet = \lim \mathcal{I}_n^\bullet$
as in Lemma \ref{lemma-K-injective}.
Then the map of the lemma becomes
$$
\lim_n 
\text{Tot}(\check{\mathcal{C}}^\bullet_{alt}(\mathcal{U},
\tau_{\geq -n}\mathcal{F}^\bullet))
\longrightarrow
\Gamma(X, \mathcal{I}^\bullet) =
\lim_n \Gamma(X, \mathcal{I}_n^\bullet)
$$
Here the arrow is in the derived category, but the equality on the
right holds on the level of complexes.
Note that (3)(b) shows that $\tau_{\geq -n}\mathcal{F}^\bullet$
is a bounded below complex satisfying the hypothesis of the lemma.
Thus the case of bounded below complexes shows each of the maps
$$
\text{Tot}(\check{\mathcal{C}}^\bullet_{alt}(\mathcal{U},
\tau_{\geq -n}\mathcal{F}^\bullet))
\longrightarrow
\Gamma(X, \mathcal{I}_n^\bullet)
$$
is a quasi-isomorphism. The cohomologies of the complexes on the left
hand side in given degree are eventually
constant (as the alternating {\v C}ech complex is finite).
Hence the same is true on the right hand side.
Thus the cohomology of the limit on the right hand side is
this constant value by
Homology, Lemma \ref{homology-lemma-apply-Mittag-Leffler-again}
(or the stronger More on Algebra, Lemma
\ref{more-algebra-lemma-apply-Mittag-Leffler-again})
and we win.
\end{proof}

\end{document}
